\documentclass[11pt]{article}
\usepackage[margin=1in]{geometry}
\usepackage{amsmath,amssymb,amsthm}
\usepackage{graphicx}
\usepackage{microtype}
\usepackage[hidelinks]{hyperref}

\newtheorem{theorem}{Theorem}
\newtheorem{lemma}{Lemma}
\newtheorem{proposition}{Proposition}
\newcommand{\R}{\mathbb{R}}
\newcommand{\Cs}{\mathcal{C}_s}
\newcommand{\M}{\mathbf{M}}
\newcommand{\supp}{\operatorname{supp}}

\title{Fractional mass does not control flat filling\\
\large A scoped counterexample for the weighted polygonal closure in arXiv:2403.03901}
\author{Automated discovery run \texttt{fa\_banach\_001}, lane 6}
\date{August 27, 2026}

\begin{document}
\maketitle

\begin{center}
\fbox{\parbox{0.91\textwidth}{\textbf{Status.}
Candidate counterexample, likely valid, pending expert review.  The result
answers the source question negatively for the closure generated by the
source paper's \emph{weighted} polygonal approximants.  A stricter closure
using only integer coefficients remains outside the claim.}}
\end{center}

\section{Source question}

The source is Cicalese--Heilmann--Kubin--Onoue--Ponsiglione,
\emph{A notion of $s$-fractional mass for $1$-currents in higher codimension},
arXiv:2403.03901v1 \cite{source}.  In the Conclusions (printed page 24), after
suggesting that ``integer'' currents might be defined as a closure of
polyhedral currents in the convergence of Theorem 5.1, the authors ask whether
a closed current of finite fractional mass bounds a 2-current of bounded local
mass, and whether the fractional mass controls the flat norm.

\begin{figure}[ht]
  \centering
  \includegraphics[width=\textwidth]{figures/open_problem_crop.png}
  \caption{The source question on printed page 24 of arXiv:2403.03901v1.}
\end{figure}

The distinction between weighted and integer coefficients matters.  Theorem
5.1 and Corollary 5.3 of the source approximate closed elements of $\Cs$ by
closed polygonal currents with coefficients $\delta_n\to0$, not by
integer-coefficient currents.  The result below treats exactly the closure
available from that approximation theorem.

\section{Definitions and statement}

Fix $d\ge3$ and $0<s<1$, and put
\[
  \beta=d-s, \qquad \theta=\frac{\beta-2}{2}
  =\frac{d-s-2}{2}>0.
\]
For a compactly supported 1-current $\gamma$ whose Fourier transform is a
function, the source norm is, up to a positive constant depending only on
$d,s$,
\[
  \|\gamma\|_{\Cs}^{2}
  \asymp \int_{\R^d}|\widehat\gamma(\xi)|^2
       |\xi|^{-\beta}\,d\xi,
  \qquad
  M_s(\gamma)=(2\pi)^{-d}\|\gamma\|_{\Cs}^{2}.
\]
A 2-current $S$ has locally finite mass if its mass measure is finite on every
compact set.  If $\partial S=\gamma$, then for every compactly supported smooth
1-form $\varphi$,
\begin{equation}\label{dual-bound}
  |\gamma(\varphi)|=|S(d\varphi)|
  \le \M(S\mathbin{\llcorner}K)
      \|d\varphi\|_{\infty},
\end{equation}
whenever $K$ contains $\supp(d\varphi)$.

\begin{theorem}[No flat filling in the weighted closure]\label{main-thm}
For every $d\ge3$ and $0<s<1$ there are a compact set $K\subset\R^d$ and a
closed 1-current $\gamma\in\Cs$, supported in $K$, with the following
properties.
\begin{enumerate}
\item There are weighted closed polygonal 1-currents $P_j$ such that
      $P_j\to\gamma$ weakly as currents and
      $M_s(P_j)\to M_s(\gamma)$.
\item There is no 2-current of locally finite mass $S$ with
      $\partial S=\gamma$.
\end{enumerate}
Consequently, the closure generated by the weighted approximation in
Theorem 5.1 and Corollary 5.3 of the source does not have the proposed filling
property, and $M_s$ cannot control its flat norm.
\end{theorem}

There is also a quantitative obstruction before taking the infinite sum.

\begin{proposition}\label{smooth-prop}
There is a sequence of smooth compactly supported divergence-free vector
fields $\gamma_N$, all supported in one compact set, such that
\[
  \sup_N M_s(\gamma_N)<\infty,
  \qquad
  \inf_{\partial S=\gamma_N}
     \M(S\mathbin{\llcorner}K)
     \longrightarrow\infty,
\]
where the infimum is over 2-currents of locally finite mass.
\end{proposition}

\section{Proof intuition}

The norm $M_s$ is a negative-order Fourier norm in dimensions $d\ge3$.
Writing a cycle as the boundary of a smooth 2-current with density
$A\sin(Nx_1)$ costs one derivative, but the source weight returns
$d-s$ derivatives in the squared norm.  The resulting $\Cs$ norm is only
$O(A N^{-\theta})$, whereas a compactly supported flat-dual test sees the
full potential amplitude $A$.  Thus a high frequency hides arbitrarily large
filling cost from fractional mass.

For the single counterexample, choose amplitudes tending to infinity and
frequencies so lacunary that the corresponding $\Cs$ norms and all unwanted
cross pairings are summable.  The cycle series converges in $\Cs$, while a
sequence of uniformly admissible flat-dual tests has pairings tending to
infinity.  Local finite mass of a filling would uniformly bound all those
pairings on their common compact support, which is impossible.

\section{Proof}

Choose nonzero real functions
$a\in C_c^\infty(\R)$ and $b\in C_c^\infty(\R^{d-1})$.  Write
$x'=(x_2,\ldots,x_d)$ and $e_{12}=e_1\wedge e_2$.  For $A>0$ and an integer
$N\ge2$, define the smooth 2-current
\[
  T_{A,N}=u_{A,N}(x)e_{12}\,dx,
  \qquad
  u_{A,N}(x)=A,a(x_1)b(x')\sin(Nx_1),
\]
and let $\gamma_{A,N}=\partial T_{A,N}$.  In vector notation,
\[
  \gamma_{A,N}=(\partial_2u_{A,N},-\partial_1u_{A,N},0,\ldots,0),
\]
so it is smooth, compactly supported, and divergence-free.

\begin{lemma}[Shift estimate]\label{shift-lemma}
There is a constant $C$, depending only on $a,b,d,s$, such that
\[
  \|\gamma_{A,N}\|_{\Cs}\le C A N^{-\theta}
  \quad\text{for all }A>0,\ N\ge2.
\]
\end{lemma}

\begin{proof}
Let $\eta(x)=a(x_1)b(x')$.  Fourier transformation gives
\[
 |\widehat{\gamma}_{A,N}(\xi)|^2
 = (\xi_1^2+\xi_2^2)|\widehat u_{A,N}(\xi)|^2,
 \qquad
 \widehat u_{A,N}(\xi)=\frac{A}{2i}
 \bigl(\widehat\eta(\xi-Ne_1)-\widehat\eta(\xi+Ne_1)\bigr).
\]
It is enough to estimate either translate.  After writing
$\zeta=\xi\mp Ne_1$, split the integral into $|\zeta|\le N/2$ and its
complement.  On the first region,
\[
 \frac{\xi_1^2+\xi_2^2}{|\xi|^\beta}
 \le |\xi|^{2-\beta}\le C N^{2-\beta}.
\]
On the complementary region, the Schwartz decay of $\widehat\eta$ gives the
same bound (indeed an error smaller than every power of $N$).  The possible
singularity at $\xi=0$ is locally integrable because
$|\xi|^{2-\beta}=|\xi|^{2-d+s}$.  Hence
\[
 \|\gamma_{A,N}\|_{\Cs}^{2}
 \le C A^2 N^{2-\beta}=C A^2N^{-2\theta},
\]
which proves the claim.
\end{proof}

\begin{lemma}[Uniform flat-dual tests]\label{dual-lemma}
There are compactly supported smooth 1-forms $\varphi_N$, all supported in one
compact set, and constants $c,C>0$ such that
\[
 \|d\varphi_N\|_\infty\le1,
 \qquad
 |\gamma_{A,N}(\varphi_N)|\ge cA
\]
for all sufficiently large $N$.  Moreover, for every $L>0$ and $m\ne n$,
\begin{multline}\label{cross-bound}
 |\gamma_{A_m,N_m}(\varphi_{N_n})|\\
 \le C_L A_m\left((1+|N_m-N_n|)^{-L}
 +(1+N_m+N_n)^{-L}\right).
\end{multline}
\end{lemma}

\begin{proof}
For a sufficiently small constant $c_0>0$, set
\[
 \varphi_N=c_0N^{-1}a(x_1)b(x')\cos(Nx_1)\,dx_2.
\]
All first derivatives of its coefficient are uniformly bounded, and choosing
$c_0$ small gives $\|d\varphi_N\|_\infty\le1$.  Since $T_{A,N}$ has only an
$e_{12}$ component,
\begin{align*}
 \gamma_{A,N}(\varphi_N)
 &=T_{A,N}(d\varphi_N)\\
 &=c_0A\int a(x_1)b(x')\sin(Nx_1)
 \left(N^{-1}a'(x_1)b(x')\cos(Nx_1)
       -a(x_1)b(x')\sin(Nx_1)\right)dx.
\end{align*}
The second term tends to
$-\frac{c_0A}{2}\|a\|_2^2\|b\|_2^2$, while the first is $o(A)$.
This proves the diagonal lower bound.  For different frequencies, the same
formula is a finite sum of Fourier coefficients of the smooth compactly
supported functions $a^2$ and $aa'$, at frequencies
$N_m-N_n$ and $N_m+N_n$.  Their rapid decay proves
\eqref{cross-bound}.
\end{proof}

Proposition \ref{smooth-prop} follows immediately: take
$A_N=N^\theta$.  Lemma \ref{shift-lemma} gives uniformly bounded fractional
mass, while Lemma \ref{dual-lemma} and \eqref{dual-bound} force every filling
to have mass at least $cN^\theta$ on the common compact support.

We now build one nonfillable current.  Put $A_j=j$.  Using Lemmas
\ref{shift-lemma} and \ref{dual-lemma}, choose integers $N_j$ inductively, as
lacunary as necessary, so that
\begin{equation}\label{lacunary-conditions}
 \|\gamma_{A_j,N_j}\|_{\Cs}\le2^{-j}
\end{equation}
and, for every $i<j$,
\begin{equation}\label{cross-small}
 |\gamma_{A_j,N_j}(\varphi_{N_i})|
 +|\gamma_{A_i,N_i}(\varphi_{N_j})|
 \le c_1 2^{-i-j-4},
\end{equation}
where $c_1>0$ is smaller than the diagonal constant in Lemma
\ref{dual-lemma}.  Such a choice is possible because $\theta>0$ and every
off-diagonal pairing tends to zero faster than any power as the new frequency
tends to infinity.

By \eqref{lacunary-conditions}, the series
\[
 \gamma=\sum_{j=1}^\infty\gamma_{A_j,N_j}
\]
converges in the weighted Fourier norm.  It also converges as a current by
weighted Cauchy--Schwarz against the Fourier transform of each test form.  All
summands have support in one compact set, so the limit is compactly supported;
its Fourier transform is the weighted $L^2$ limit, hence $\gamma\in\Cs$.
Also $\partial\gamma=0$.

The diagonal estimate and \eqref{cross-small} give
\[
 |\gamma(\varphi_{N_j})|
 \ge c_1A_j-\sum_{m\ne j}
      |\gamma_{A_m,N_m}(\varphi_{N_j})|
 \ge \frac{c_1}{2}j
\]
after decreasing $c_1$ if necessary.  Thus these pairings are unbounded,
although $\|d\varphi_{N_j}\|_\infty\le1$ and all $d\varphi_{N_j}$ have common
compact support.  Equation \eqref{dual-bound} rules out every locally
finite-mass 2-current $S$ with $\partial S=\gamma$.

Finally, Corollary 5.3 of the source applies to every closed
$\gamma\in\Cs$ and supplies weighted closed polygonal currents converging
weakly to $\gamma$ with convergence of $M_s$.  This proves both assertions of
Theorem \ref{main-thm}. \qed

\section{Verification notes and limitations}

The proof is symbolic.  Its two threshold checks are sharp for this method:
$\theta>0$ exactly when $d-s>2$, and the source explicitly records the planar
case as positive.  The singular part of the Fourier estimate is integrable,
and translated Schwartz decay controls the only region where the simple
$N^{2-\beta}$ pointwise estimate fails.

The conclusion is stronger than failure of a uniform bound: the lacunary
limit has no filling of locally finite mass at all.  On the other hand, the
coefficient caveat is essential.  The proof uses the weighted closure provided
by Corollary 5.3.  It does not prove that $\gamma$ is a limit of
integer-coefficient polyhedral currents in the same topology.  Therefore a
classically integral-coefficient version of the source question remains open.

Human review should focus on Lemma \ref{shift-lemma}, the simultaneous
lacunary choice in \eqref{lacunary-conditions}--\eqref{cross-small}, and the
intended meaning of ``integer'' in the source Conclusions.

\section{Novelty search}

On August 27, 2026, a bounded search covered the run's registry, solution,
attempt, and proof-gap indexes; arXiv:2403.03901; the exact title and author
combination; and arXiv searches combining \texttt{fractional mass},
\texttt{flat norm}, \texttt{weighted polygonal}, \texttt{Plateau}, and
\texttt{1-currents}.  The arXiv record has only v1.  No later primary-source
answer or the oscillatory/lacunary obstruction above was found.  This is not a
claim of exhaustive priority.

\begin{thebibliography}{9}
\bibitem{source}
M. Cicalese, T. Heilmann, A. Kubin, F. Onoue, and M. Ponsiglione,
\emph{A notion of $s$-fractional mass for $1$-currents in higher codimension},
arXiv:2403.03901v1, 2024.
\end{thebibliography}

\end{document}
